\section{Arximed qonuni}


\begin{enumerate}

\item Suvda to'liq cho'kkan, 0,5 $m^3$ hajmli toshga ta'sir etadigan itarib chiquvchi kuchni (kN da) aniqlang. $g=10m/s^2$.
\begin{enumerate}[label=\Alph*)]
\item 10
\item 5
\item 15
\item 2,5
\end{enumerate}

\item Havoda 50 $dm^3$ hajmli jismga ta'sir qiladigan arximed kuchi (mN da) qanchaga teng? Havoning zichligi 1,29 $kg/m^3$. $g=10m/s^2$.
\begin{enumerate}[label=\Alph*)]
\item 645
\item 500
\item 129
\item 322
\end{enumerate}

\item Havodagi og'irligi 350 N bo'lgan toshni suv ostida ushlab turish uchun qanday kuch qo'yish kerak? Tosh moddasining zichligi 2500 $kg/m^3$.
\begin{enumerate}[label=\Alph*)]
\item 140
\item 280
\item 210
\item 175
\end{enumerate}

\item Massasi 2 kg va hajmi 0,001 $m^3$ bo'lgan jism ko'l ichida 5 m chuqurlikda joylashgan. Uni suv sirtidan 5 m balandlikkacha sekinlik bilan ko'tarish uchun qanday ish bajarish kerak? $g=10m/s^2$.
\begin{enumerate}[label=\Alph*)]
\item 100
\item 200
\item 75
\item 150
\end{enumerate}

\item Alyuminiy bo'lagining havodagi og'irligi 270 N, glitserindagi og'irligi esa 144 N. Agar alyuminiyning zichligi 2700 $kg/m^3$ bo'lsa, glitserinning zichligini aniqlang.
\begin{enumerate}[label=\Alph*)]
\item 1260
\item 1000
\item 800
\item 1150
\end{enumerate}

\item Jismning suvdagi og'irligi havodagi og'irligidan 2 marta kichik. Jism moddasining zichligi ($g/cm^3$ da) qanday?
\begin{enumerate}[label=\Alph*)]
\item 1,5
\item 2,5
\item 2
\item 3
\end{enumerate}

\item Agar qandaydir bir jismning zichligi 800 $kg/m^3$ bo'lgan suyuqlikdagi og'irligi havodagi og'irligiga qaraganda 3 marta kichik bo'lsa, jismning zichligi shu suyuqlik zichligidan qanchaga katta?
\begin{enumerate}[label=\Alph*)]
\item 400
\item 800
\item 1200
\item 600
\end{enumerate}

\item Bir jismning zichligi suvning zichligidan 1,25 marta katta. Shu jismning suvdagi og'irligi havodagi og'irligidan necha marta kichik bo'ladi?
\begin{enumerate}[label=\Alph*)]
\item 4
\item 5
\item 1,25
\item 2,5
\end{enumerate}

\item Metall bo'lagining suvdagi og'irligi 6800 N, benzindagi og'irligi esa 7100 N. Benzinning zichligini aniqlang. Metallning zichligi 7800 $kg/m^3$.
\begin{enumerate}[label=\Alph*)]
\item 800
\item 900
\item 700
\item 750
\end{enumerate}

\item Jism havoda 3 N, suvda 1,8 N va noma'lum zichlikdagi suyuqlikda esa 2,04 N og'irlikka ega. Shu noma'lum suyuqlikning zichligi qanday?
\begin{enumerate}[label=\Alph*)]
\item 600
\item 800
\item 1000
\item 900
\end{enumerate}

\item Probka bo'lagi havoda 1 N og'irlikka, metall bo'lagi esa 10 N og'irlikka ega. Agar bu bo'laklarni ip bilan bog'lab kerosinga to'liq botirilsa, ularning umumiy og'irligi 5 N bo'ladi. Probkaning zichligini toping. Kerosinning zichligi 800 $kg/m^3$, metallning zichligi 4000 $kg/m^3$.
\begin{enumerate}[label=\Alph*)]
\item 200
\item 400
\item 150
\item 300
\end{enumerate}

\item Kvars ichiga joylangan sof oltin u bilan birgalikda dinamometr prujinasini 2,26 N kuch bilan cho'zadi. Oltin suvga botirilganda unga 0,2 N itarib chiquvchi kuch ta'sir qiladi. Oltinning massasini (g da) toping. Oltinning zichligi 19,3 $g/cm^3$, kvarsning zichligi 3,3 $g/cm^3$. $g=10m/s^2$.
\begin{enumerate}[label=\Alph*)]
\item 100
\item 193
\item 226
\item 150
\end{enumerate}

\item Oltin va kumush qotishmasidan tayyorlangan toj tarkibidagi oltinning massasini (g da) toping. Tojning havodagi og'irligi 25,4 N, suvdagisi esa 23,4 N. Oltinning zichligi 19,3 $g/cm^3$, kumushniki - 10,5 $g/cm^3$. $g=10m/s^2$.
\begin{enumerate}[label=\Alph*)]
\item 500
\item 800
\item 965
\item 1200
\end{enumerate}

\item 1 $cm^3$ hajmli shisha sharcha suvda tekis tushib bormoqda. Sharchaning 10 m ga ko'chishida 0,17 J issiqlik ajraladi. Shishaning zichligini toping. $g=10m/s^2$.
\begin{enumerate}[label=\Alph*)]
\item 2500
\item 2700
\item 2000
\item 3000
\end{enumerate}

\item Suvli idish ichida shar shunday yotibdiki, bunda uning yarmi suvda joylashgan va u idish tubiga tegib turadi. Agar sharning havodagi og'irligi 8 N bo'lsa, u idish tubiga qanday kuch bilan bosadi? Shar moddasining zichligi 800 $kg/m^3$.
\begin{enumerate}[label=\Alph*)]
\item 5
\item 3
\item 4
\item 2
\end{enumerate}

\item 60 g massali bir jinsli sharcha bo'sh stakanning ostida yotibdi. Stakanga suyuqlik shunday quyiladiki, bunda sharcha hajmining suyuqlikka botgan qismi uning umumiy hajmidan 6 marta kichik bo'lib qoladi. Suyuqlikning zichligi sharcha materialining zichligidan 3 marta katta. Sharchaning stakan tubiga bosim kuchini (mN da) toping. $g=10m/s^2$.
\begin{enumerate}[label=\Alph*)]
\item 300
\item 600
\item 150
\item 450
\end{enumerate}

\item Uzunligi 2 m va ko'ndalang kesim yuzasi 100 $cm^2$ bo'lgan alyuminiyli silindr suv havzasidan o'zgarmas tezlik bilan asta-sekin tortib chiqariladi. Havzadan silindr uzunligining 0,25 ga teng qismi chiqqan paytda arqon uzilib ketdi. Arqonning maksimal tarangligini aniqlang. Alyuminiyning zichligi 2700 $kg/m^3$, $g=10m/s^2$.
\begin{enumerate}[label=\Alph*)]
\item 540
\item 390
\item 450
\item 300
\end{enumerate}

\item Ipning bir uchi tubga, ikkinchisi esa probkali po'kakka mahkamlangan. Bunda po'kakning butun hajmining 75\% i suvga botgan. Agar po'kakning massasi 2 kg ga teng bo'lsa, ipning taranglik kuchini aniqlang. Po'kakning zichligi 300 $kg/m^3$, $g=10m/s^2$.
\begin{enumerate}[label=\Alph*)]
\item 20
\item 40
\item 30
\item 50
\end{enumerate}

\item Ichida suvi bor silindrik idishda muz parchasi ip bilan idish tubiga tortilgan holda suzib yuribdi. Muz parchasi erigach, suv sathi 1 cm ga pasaydi. Ipning taranglik kuchi qanday bo'lgan? Idish tubining yuzasi 100 $cm^2$. $g=10m/s^2$.
\begin{enumerate}[label=\Alph*)]
\item 1
\item 10
\item 0,1
\item 5
\end{enumerate}

\item Yuzi 300 $cm^2$ bo'lgan suvli silindrik idishning tubida 3000 $kg/m^3$ zichlikdagi materialdan tayyorlangan, asosining yuzi 100 $cm^2$ va balandligi 40 cm bo'lgan silindr turibdi. Agar suv qatlamining boshlang'ich qalinligi 30 cm bo'lsa, silindrni suvdan chiqarib olish uchun qanday ish bajarish kerak? $g=10m/s^2$.
\begin{enumerate}[label=\Alph*)]
\item 21
\item 42
\item 10
\item 30
\end{enumerate}

\item Yuzi 400 $cm^2$ bo'lgan suvli silindrik idishning tubida 2500 $kg/m^3$ zichlikdagi materialdan tayyorlangan, asosining yuzi 100 $cm^2$ va balandligi 40 cm bo'lgan silindr turibdi. Agar suv qatlamining boshlang'ich qalinligi 60 cm bo'lsa, silindrni suvdan chiqarib olish uchun qanday ish bajarish kerak? $g=10m/s^2$.
\begin{enumerate}[label=\Alph*)]
\item 20
\item 40
\item 36
\item 54
\end{enumerate}

\item Yupqa bir jinsli tayoqcha yuqori uchidan sharnirli mahkamlangan. Tayoqchaning pastki qismi suvga botirilgan. Tayoqcha suv yuzasiga nisbatan qiya joylashib, uning yarmi suv ichida bo'lgan vaqtda muvozanat vaziyati ro'y beradi. Tayoqcha yasalgan materialning zichligi qanday?
\begin{enumerate}[label=\Alph*)]
\item 500
\item 750
\item 800
\item 600
\end{enumerate}

\item Balandligi 8 cm va radiusi 3 cm bo'lgan silliq stakan ichiga 12 cm uzunlikdagi va 100 g massali bir jinsli tayoqcha qo'yildi. Stakan zichligi tayoqcha materialining zichligidan ikki marta kichik bo'lgan suyuqlik bilan to'ldirildi. Tayoqcha stakanning chetini qanday kuch (mN da) bilan bosadi? $g=10m/s^2$.
\begin{enumerate}[label=\Alph*)]
\item 235
\item 470
\item 117
\item 300
\end{enumerate}

\item 4 cm radiusli silliq baland stakan ichiga uzunligi 10 cm va massasi 60 g bo'lgan tayoqcha qo'yildi, so'ngra stakanga 3 cm balandlikkacha zichligi tayoqcha materialining zichligidan bir yarim marta katta bo'lgan suyuqlik quyildi. Tayoqchaning yuqori uchi stakanning devoriga qanday kuch bilan (mN da) bosadi? $g=10m/s^2$.
\begin{enumerate}[label=\Alph*)]
\item 125
\item 500
\item 250
\item 375
\end{enumerate}

\item Yaxlit jism suvda suzmoqda, bunda jism hajmining 3/4 qismi suvga botgan. Jismning hajmi 0,1 $m^3$. Jismga ta'sir etuvchi og'irlik kuchini aniqlang. $g=10m/s^2$.
\begin{enumerate}[label=\Alph*)]
\item 1000
\item 500
\item 750
\item 250
\end{enumerate}

\item Yog'och bo'lagi hajmining 3/4 qismi suvga botgan holda suvda suzmoqda. Yog'och qanday zichlikka ega?
\begin{enumerate}[label=\Alph*)]
\item 750
\item 800
\item 600
\item 500
\end{enumerate}

\item Hajmi 0,0002 $m^3$ bo'lgan bir jinsli jism zichligi jism materiali zichligidan 4 marta katta bo'lgan suyuqlikda suzmoqda. Jismning qanday hajmi ($cm^3$ da) suyuqlik sirtidan chiqib turadi?
\begin{enumerate}[label=\Alph*)]
\item 50
\item 100
\item 150
\item 200
\end{enumerate}

\item Zichligi 800 $kg/m^3$ bo'lgan jism zichligi 1000 $kg/m^3$ bo'lgan suyuqlikda suzadi. Jism hajmining qanday qismi (foizlarda) suyuqlikka botgan?
\begin{enumerate}[label=\Alph*)]
\item 20
\item 80
\item 50
\item 75
\end{enumerate}

\item Suvda suzuvchi suv ichida chalqancha holda harakatsiz yotibdi. Bunda uning yuzining oz qismidan tashqari, butun tanasi suvga botgan. Suzuvchining massasi 75 kg. Uning tanasining hajmini ($dm^3$ da) toping.
\begin{enumerate}[label=\Alph*)]
\item 75
\item 80
\item 70
\item 100
\end{enumerate}

\item Kemaga ta'sir qiluvchi og'irlik kuchi $10^3$ kN ga teng. Bu kema qanday hajmdagi suvni siqib chiqaradi?
\begin{enumerate}[label=\Alph*)]
\item 10
\item 1000
\item 100
\item 50
\end{enumerate}

\item Tubining yuzasi 38 $cm^2$ va balandligi 6 cm bo'lgan to'g'ri burchakli tunuka quticha suvda shunday suzmoqdaki, bunda uning suv yuzidagi qismining balandligi 4 cm ga teng. Qutichaning massasini (g da) aniqlang.
\begin{enumerate}[label=\Alph*)]
\item 38
\item 76
\item 152
\item 114
\end{enumerate}

\end{enumerate}

\begin{table}[h!]
\centering
\caption{Arximed qonuni va Suzish shartlari uchun to'g'ri javoblar jadvali (1-31)}
\begin{tabular}{|c|c|c|c|c|c|c|c|c|c|c|c|}
\hline
№ & J. & № & J. & № & J. & № & J. & № & J. & № & J. \\ \hline
1 & B & 7 & A & 13 & C & 19 & A & 25 & C & 31 & B \\ \hline
2 & A & 8 & B & 14 & B & 20 & A & 26 & A & & \\ \hline
3 & C & 9 & C & 15 & B & 21 & C & 27 & C & & \\ \hline
4 & D & 10 & B & 16 & A & 22 & B & 28 & B & & \\ \hline
5 & A & 11 & A & 17 & B & 23 & A & 29 & A & & \\ \hline
6 & C & 12 & B & 18 & C & 24 & C & 30 & C & & \\ \hline
\end{tabular}
\end{table}